\documentclass[10pt,journal,compsoc]{IEEEtran}

\ifCLASSOPTIONcompsoc
  \usepackage[nocompress]{cite}
\else
  \usepackage{cite}
\fi

\usepackage{amsmath}
\usepackage{graphicx}
\usepackage{url}

\hyphenation{op-tical net-works semi-conduc-tor}

\begin{document}

\title{JOURNAL}

\author{
\begin{tabular}{ccc}
\textbf{Arka Dutta} &
\textbf{Naveed Wazeed Khan} &
\textbf{Hasan Murad} \\

\texttt{u2204025@student.cuet.ac.bd} &
\texttt{u2204016@student.cuet.ac.bd} &
\texttt{hasanmurad@cuet.ac.bd}
\end{tabular}
}

\markboth{IEEE Journal Submission Draft}%
{Author \MakeLowercase{\textit{et al.}}: Bengali Government-Service Question Answering}

\IEEEtitleabstractindextext{%
\begin{abstract}
% Write your abstract here.
\end{abstract}

\begin{IEEEkeywords}
% Add IEEE index terms / keywords here.
\end{IEEEkeywords}}

\maketitle

\IEEEdisplaynontitleabstractindextext
\IEEEpeerreviewmaketitle
\section{Introduction}
\label{sec:introduction}

The growing use of digital technologies in government services has changed how citizens obtain administrative information and communicate with public institutions. Information about services such as birth and death registration, National Identification (NID), passport and e-passport applications, and Taxpayer Identification Number (TIN) registration is now widely available online. However, users may still need to navigate multiple webpages to find information about procedures, required documents, fees, eligibility criteria, processing times, and other service requirements. Conversational artificial intelligence (AI) can make this information easier to access by allowing users to ask questions in everyday language rather than searching through several government websites. For this reason, chatbots and conversational systems are becoming increasingly relevant to digital government and public-service delivery~\cite{cortes2023egov}.

Recent developments in large language models (LLMs) have improved the ability of conversational systems to understand natural-language queries and generate relevant responses. Modern generative models can perform tasks such as question answering, summarization, instruction following, and multilingual generation with relatively little task-specific engineering. Using these models for government-service question answering, however, raises concerns about the reliability of their responses. Since LLMs mainly depend on knowledge acquired during training, they may provide information that appears reliable but is incorrect, incomplete, outdated, or unsupported. In public-service applications, such errors can have practical consequences for citizens.

These challenges become more significant when services are provided in Bangla. Bangla has a large number of speakers, but it is still less represented in NLP resources and pretrained language models. Recent studies have explored Bangla-specific language models and instruction datasets to address this gap. For instance, BanglaLlama demonstrates that instruction tuning with Bangla data can improve LLM performance on Bangla-language tasks~\cite {zehady-etal-2026-banglallama}.

Government-service question answering requires more than general Bangla language understanding. User queries may include Bangla script, English administrative terms, numerical information, dates, fees, document names, and URLs. An effective system should be able to handle these different language patterns while providing accurate answers based on reliable government information. Several strategies can be used to adapt general-purpose LLMs to this domain. Prompt-based approaches, including zero-shot, few-shot, and chain-of-thought prompting, can guide model responses without changing model parameters. Supervised fine-tuning can further adapt a model to the language patterns, terminology, and response structure of government-service questions. Fine-tuning alone does not fully address the problem of knowledge freshness. Government procedures, fees, eligibility requirements, and online resources can change over time, while updating information stored in model parameters generally requires additional training. A system that relies only on parametric knowledge may therefore become outdated even when its language-generation capabilities remain strong.

Retrieval-Augmented Generation (RAG) offers a practical way to address this limitation. Rather than relying solely on information stored in model parameters, a RAG system retrieves relevant documents or passages from an external knowledge source and provides them to the language model as context before generating an answer~\cite{lewis2020rag}. This approach is particularly suitable for government-service applications because the knowledge base can be updated independently of the generative model. When a government procedure, fee, requirement, or online resource changes, the external knowledge source can be updated without retraining the entire LLM. Retrieval can also connect generated answers to their supporting sources, which can improve transparency and trust in citizen-facing systems.

This study investigates different LLM adaptation strategies for Bangla government-service question answering. For this purpose, we developed a domain-specific question--answer dataset covering four government-service categories in Bangladesh: birth and death registration, passport and e-passport services, National Identity services, and TIN-related services. The dataset contains 1,186 question--answer instances, including 938 training samples and 248 test samples. We also prepared domain-specific knowledge bases to support retrieval-based question answering. We evaluated different model families, including Llama, Qwen, Mistral, and Bangla-focused language models. The models were evaluated under different configurations, including zero-shot, few-shot, and chain-of-thought prompting, supervised fine-tuning, RAG, and a combination of RAG and fine-tuning.
The main aim of this study is to examine whether using external government-service information improves answer quality compared with prompting and supervised fine-tuning alone. We also investigate whether combining retrieval with domain-specific fine-tuning provides further improvements. To compare these approaches, we use several evaluation metrics, including Exact Match, fuzzy matching, BLEU, ROUGE, Token F1, and BERTScore. The experimental results show that prompting-only methods generally perform worse than retrieval-based approaches for government-service question answering. Among the evaluated methods, RAG combined with supervised fine-tuning gives the best overall results. This shows that using retrieved information together with domain-specific training can help improve the quality of answers.

The main contributions of this work are as follows:
\begin{itemize}
    \item We focus on Bangla question answering for four government-service domains in Bangladesh: birth and death registration, passport and e-passport services, National Identity services, and TIN services.

    \item We develop a domain-specific question--answer dataset and prepare knowledge bases that support both supervised fine-tuning and RAG.

    \item We compare different approaches, including zero-shot, few-shot, and chain-of-thought prompting, supervised fine-tuning, RAG, and RAG combined with fine-tuning across different model families.

    \item We identify the limitations of traditional evaluation metrics and emphasize the need to verify answers using reliable and up-to-date government sources.
\end{itemize}
The rest of the paper is organized as follows. Section~\ref{sec:Related_work} reviews previous work on Bangla language models, government-service conversational systems, LLM adaptation, and retrieval-augmented generation. Section~\ref{sec:Dataset and Knowledge Base} describes the dataset and the process used to build the knowledge base. Section~\ref{sec:Proposed Methodology} presents the methodology and experimental configurations. Section~\ref{sec:Experimental Setup} details the experimental setup and evaluation metrics, while Section~\ref{sec:Results And Dis} presents and discusses the results. Section~\ref{sec:limitations} covers the limitations, reliability concerns, and practical aspects of deployment. Finally, Section~\ref{sec:conclusion} concludes the paper with a summary of the main findings and possible directions for future research.


\section{Related Work}
\label{sec:Related_work}
       The growing use of large language models (LLMs) is creating new ways to provide public information. They can make access to government services easier through conversational systems. In e-government applications, these models can assist users by making complex administrative information easier to access and understand through natural language interactions. However, their effective use in real-world public-service settings largely depends on their ability to provide responses that are accurate, relevant, and grounded in reliable domain-specific information.Research on LLM-based digital-government systems has started to address these challenges.
       
     Gao \textit{et al.} \cite{gao2023digitalgovernment} investigated the use of LLMs for intelligent question answering in digital-government applications, focusing on user-intent recognition, conversational context, and complex questions. Their work demonstrates the potential of combining the contextual capabilities of LLMs with conventional information-retrieval mechanisms. Similarly, Mamalis \textit{et al.} \cite{mamalis2024opengovernment} explored the integration of GPT-3.5 with the Scottish open-statistics portal. Their study was motivated by the tendency of standalone LLMs to generate factually incorrect information and demonstrated how authoritative government data can be connected to an LLM-based natural-language interface. These studies indicate that the usefulness of LLMs in public-service applications depends not only on fluent response generation but also on their ability to access reliable domain-specific information.

The need for external grounding is closely related to the limitations of parametric knowledge in language models. Roberts \textit{et al.} \cite{roberts2020knowledge} investigated how much factual knowledge could be stored within the parameters of pretrained language models. Their findings showed that larger models can answer a considerable number of factual questions using internalized knowledge alone. However, such closed-book approaches depend entirely on information acquired during pretraining and do not provide direct access to updated external evidence. This limitation is particularly important for government-service applications, where fees, application procedures, eligibility requirements, required documents, and other administrative information may change over time.

Retrieval-Augmented Generation (RAG) provides an important approach to addressing this limitation. Lewis \textit{et al.} \cite{lewis2020rag} introduced the RAG framework by combining parametric knowledge stored in a pretrained sequence-to-sequence model with non-parametric knowledge retrieved from an external document collection. Rather than requiring the generator to rely solely on information encoded within its parameters, a retrieval component first identifies relevant documents or passages and provides them as contextual evidence during response generation. Their results demonstrated the effectiveness of retrieval augmentation for knowledge-intensive natural-language processing tasks and established RAG as a foundation for systems requiring access to external and updatable knowledge.

Several recent studies have applied retrieval-augmented approaches specifically to government and public-service environments. Pujiono  \textit{et al.}  \cite{pujiono2024rag} looked at how to combine RAG, vector databases, and large language models for chatbot use in public-service agencies. Their method first retrieves relevant regulatory and service-related information before creating a response. This helps make the answers more accurate and based on data.

Papageorgiou \textit{et al.} \cite{papageorgiou2024egovernment} proposed a modular and reproducible LLM-based architecture for e-government services that combines RAG with external knowledge sources. The architecture uses structured retrieval components to support scalability and improve the reliability of government-service conversational systems.


More recently, Jiang \textit{et al.} \cite{jiang2025argus} introduced ARGUS, a retrieval-augmented question-answering system for government services. It uses domain adaptation and hybrid retrieval to improve the factual grounding of LLM-generated responses. Their approach combines multiple retrieval techniques rather than relying on a single vector-search mechanism and evaluates the resulting system using both conventional question-answering metrics and retrieval-oriented evaluation measures. This work highlights that the performance of a government-service chatbot depends not only on the underlying language model but also on the quality of the retrieval mechanism and the domain knowledge supplied to the model.

In addition to retrieval, supervised fine-tuning has been widely studied as a method for adapting general-purpose LLMs to specific tasks and domains. Fine-tuning enables a model to learn domain-specific vocabulary, response patterns, and task behavior from specialized training examples. However, retrieval and fine-tuning address different aspects of model adaptation. Ovadia \textit{et al.} \cite{ovadia2024finetuning} compared fine-tuning and retrieval-based approaches for injecting knowledge into LLMs. Their experiments showed that retrieval augmentation can provide substantial advantages for knowledge-intensive tasks, particularly when models need access to newly introduced information that may not have been represented during training.

Alghisi \textit{et al.} \cite{alghisi2024finetunerag} further compared several LLM adaptation strategies, including in-context learning, fine-tuning, RAG, and knowledge grounding. Their experiments, involving models such as Llama-2 and Mistral, showed that the effectiveness of each adaptation method varies depending on the model and the task. Consequently, no single strategy can be assumed to perform best across all conversational settings. These findings support fair comparisons of adaptation strategies under the same experimental conditions.

Recent studies have also examined the use of retrieval together with fine-tuning. Lawton \textit{et al.} \cite{lawton2025ragfinetuning} explored independent, joint, and multi-stage fine-tuning approaches within RAG pipelines. Their findings show that retrieval and fine-tuning can be combined through different strategies. These strategies differ in both computational requirements and model performance. This is especially relevant for domain-specific question answering because fine-tuning can improve task-specific response behavior, while retrieval provides access to external factual information. Therefore, RAG and fine-tuning need not be considered mutually exclusive strategies.

Another important consideration is the choice of the underlying language model. Recent open-weight model families, including Mistral, Qwen, and Llama, provide strong alternatives for developing domain-specific conversational systems. Jiang \textit{et al.} \cite{jiang2023mistral} introduced Mistral 7B, an efficient seven-billion-parameter model incorporating grouped-query attention and sliding-window attention. Their experiments demonstrated that relatively compact models can achieve competitive performance compared with substantially larger architectures. Yang \textit{et al.} \cite{yang2024qwen2} introduced Qwen2, a multilingual LLM family covering different parameter scales and supporting language understanding, generation, reasoning, mathematics, coding, and instruction-following tasks. Its multilingual capabilities make the model particularly relevant to applications involving languages other than English.

Grattafiori \textit{et al.} \cite{grattafiori2024llama3} introduced the Llama 3 family of open-weight foundation models for general-purpose language understanding, multilingual generation, reasoning, coding, and instruction following. Llama models have subsequently become widely used as open-weight baselines for downstream LLM applications. Collectively, Qwen, Mistral, and Llama provide suitable model families for controlled comparisons of domain-specific question-answering systems. However, their original evaluations primarily focus on general-purpose benchmarks rather than Bangla government-service question answering.

The performance of these models on Bangla cannot be inferred directly from their results on English and other high-resource languages. Bangla remains comparatively under-resourced in terms of high-quality question-answering datasets and comprehensive benchmarks. Aurpa \textit{et al.} \cite{aurpa2022banglaqa} investigated reading-comprehension-based question answering in Bangla using recurrent and transformer-based architectures. Their results demonstrated considerable improvements from transformer-based approaches over recurrent models, illustrating the potential of pretrained transformers for Bangla question answering. Nevertheless, their study primarily focused on reading comprehension rather than generative and retrieval-augmented question answering.

Ekram \textit{et al.} \cite{ekram2022banglarqa} introduced BanglaRQA, a Bangla reading-comprehension benchmark developed to address the limited availability of comprehensive resources for Bangla question answering. Their experiments with transformer-based models demonstrated substantial differences in performance across models and question categories. These findings suggest that aggregate scores may not fully reflect model performance. Performance can vary across different types of Bangla questions.

Kabir  \textit{et al.} \cite{kabir2024benllm} examined the limitations of general-purpose LLMs for Bengali NLP using the BenLLM-Eval benchmark. The study evaluated several LLMs on tasks such as question answering, summarization, paraphrasing, classification, natural language inference, and sentiment analysis. The results showed clear differences across models and tasks. A model that performs well in general settings may not perform equally well on Bengali NLP tasks. This finding supports the need for task-specific evaluation when applying LLMs to Bangla government-service applications.

Chang  \textit{et al.} \cite{chang2024llmevaluation}  provided a broad survey of LLM evaluation. They showed that model performance should be assessed from several perspectives instead of using only one metric. Common metrics such as Exact Match, F1-score, BLEU, and ROUGE measure different forms of text similarity. However, they do not always show whether an answer is factually correct or supported by the retrieved evidence. This is especially important in government-service applications. A response may look similar to the reference answer but still contain wrong information about fees, procedures, required documents, or eligibility rules.
This limitation becomes particularly important in government-service applications, where an answer may be linguistically similar to a reference while still containing an incorrect fee, procedure, required document, or eligibility condition.
Existing literature therefore points to three connected research directions. One important direction in e-government research is ensuring that responses generated by language models are grounded in reliable and authoritative information sources. Another direction concerns the comparison of retrieval-based and fine-tuning approaches, with existing studies showing that their effectiveness depends on the specific task and the model being used. A further direction is evident in Bangla question-answering research, where transformer-based models have shown promising progress while challenges remain in low-resource language understanding and performance evaluation.

Despite these developments, limited research has jointly investigated these issues in a controlled Bangla government-service setting. Existing government-oriented LLM studies have primarily considered English or other comparatively high-resource languages, specific system architectures, or individual retrieval strategies. In contrast, existing Bangla question-answering research has largely concentrated on reading comprehension and general Bengali NLP benchmarks rather than retrieval-grounded government-service conversations. Furthermore, comparative studies of prompting, supervised fine-tuning, RAG, and their combination remain limited for Bangla administrative information.

These limitations motivate a controlled comparison of prompting, supervised fine-tuning, RAG, and hybrid adaptation for Bangla government-service question answering. Unlike earlier work, this study evaluates these strategies across multiple open-weight model families using the same domain-specific dataset and evaluation setting.
\section{Dataset}
\label{sec:dataset}
\label{sec:Dataset and Knowledge Base}

We constructed a Bangla question answering dataset and a new knowledge base covering four public service domains in Bangladesh: passport and e-passport services, National Identity (NID) services, birth and death registration, and Taxpayer Identification Number (TIN) services. Information was collected by crawling official government websites~\cite{gov_epassport,gov_nid,gov_bdris,gov_tin}. The knowledge base was prepared first, and question and answer pairs were then generated from its contents. This provided a common source of information for studying prompting, fine-tuning, retrieval-augmented generation (RAG), and their combination.

The selected domains cover several kinds of administrative questions. Passport queries concern applications, renewals, supporting documents, fees, and delivery periods. NID queries include registration, correction of personal information, replacement, and smart card services. Birth and death registration form the civil registration domain, which contains information about registration procedures, certificates, and corrections. The TIN domain covers taxpayer registration and related services. Across these domains, an answer may require a short factual statement, a list of documents, or a sequence of application steps.

The collected webpages contained service information alongside navigation menus and other material that was not needed for question answering. We processed the extracted content to retain the relevant information and used GPT-5~\cite{openai2025gpt5} to assist in organizing it by domain. The resulting knowledge base was stored as an array of JSON objects. Records included the service information, source details, and keywords, with metadata such as document identifier, domain, topic, and title. Retaining the source URL allowed the information in a record to be traced to its government website.

% Figure content: Data preparation workflow. Show official government websites, crawling and content preparation, the four domain knowledge bases, QA generation, and the training and test partitions. Show a separate branch from the knowledge base to the retrieval index.
\begin{figure}[!t]
\centering
% Replace the box below with \includegraphics[width=\columnwidth]{your_figure_file}.
\fbox{\parbox[c][1.15in][c]{0.94\columnwidth}{\centering Figure placeholder}}
\caption{Construction of the government service knowledge base and question answering dataset.}
\label{fig:data_preparation}
\end{figure}

After preparing the knowledge base, we prompted GPT-5 to generate questions and reference answers from the collected information. The instructions requested coverage of the available service details, including eligibility, required documents, application procedures, fees, corrections, and processing times. We also requested paraphrased questions to represent different ways of asking for the same information. The answers were generated from the knowledge base. Thus, the dataset consists of questions prepared for this study, rather than queries collected from citizen interactions.

Each QA record contains an instruction field for the question, an input field for additional context where available, and an output field for the reference answer. Supporting fields identify the domain and source, with further metadata where provided. The questions include Bangla script as well as variations containing English administrative terms and romanized Bangla. This variation reflects the language forms represented in the prepared dataset. Figure~\ref{fig:dataset_examples} illustrates examples from the four domains and their corresponding knowledge base records.

% Figure content: Examples taken directly from the stored dataset. Include one question and reference answer from each domain, with the corresponding source URL. Alongside one example, show the structure of its knowledge base record. Preserve the original wording.
\begin{figure}[!t]
\centering
% Replace the box below with \includegraphics[width=\columnwidth]{your_figure_file}.
\fbox{\parbox[c][1.15in][c]{0.94\columnwidth}{\centering Figure placeholder}}
\caption{Examples and record structure of the Bangla government service dataset.}
\label{fig:dataset_examples}
\end{figure}

The final dataset contains 1,186 question and answer pairs. Of these, 938 were assigned to training and 248 to testing. Table~\ref{tab:dataset_summary} presents the distribution by service domain. Birth and death registration accounts for 330 pairs, followed by passport services with 312, NID services with 307, and TIN services with 237. These figures describe the prepared dataset; they do not represent the frequency of citizen requests for each service.

\begin{table}[!t]
\centering
\caption{Distribution of question and answer pairs across the training and test sets.}
\label{tab:dataset_summary}
\footnotesize
\setlength{\tabcolsep}{3pt}
\renewcommand{\arraystretch}{1.1}
\begin{tabular}{lrrr}
\hline
\textbf{Domain} & \textbf{Train} & \textbf{Test} & \textbf{Total} \\
\hline
Birth and death registration & 267 & 63 & 330 \\
Passport and e-passport & 249 & 63 & 312 \\
NID & 233 & 74 & 307 \\
TIN & 189 & 48 & 237 \\
\hline
\textbf{Total} & \textbf{938} & \textbf{248} & \textbf{1,186} \\
\hline
\end{tabular}
\end{table}

The resources served different purposes in the experiments. Fine-tuning used the 938 training pairs to adapt the generators to the task. Prompting experiments used the test questions, with training examples supplied as demonstrations where applicable. RAG retrieved passages from the government service knowledge base, while the combined approach used those passages with a generator adapted on the training pairs. All configurations were evaluated on the same 248 test questions, and their reference answers were withheld from the inference prompts.

The knowledge base remained available to the retrieval systems during evaluation. Since the QA pairs were generated from this collection, the test set measures performance on held-out questions about the collected information. The following section describes how the models and retrieval components were configured for these comparisons.

\section{Experiments}
\label{sec:experiments}
% Compatibility labels preserve existing introduction cross-references.
\label{sec:Proposed Methodology}
\label{sec:Experimental Setup}

We evaluated prompting, supervised fine-tuning, retrieval-augmented generation (RAG), and RAG combined with fine-tuning for Bangla government service question answering. The experiments used Qwen2.5-7B-Instruct~\cite{yang2024qwen25}, Llama-3.1-8B-Instruct~\cite{grattafiori2024llama3}, and Mistral-7B-Instruct-v0.3~\cite{jiang2023mistral} as answer generators. Each model was assessed on the same 248 test questions described in Section~\ref{sec:dataset}. The comparison examined whether task-specific training, retrieved information, or their combination improved agreement with the reference answers.

All experiments were carried out on Kaggle using an NVIDIA T4 GPU. Prompting and retrieval-based inference used 4-bit model representations. The pretrained generators were loaded with NF4 quantization, double quantization, and FP16 computation~\cite{meth_dettmers2023qlora}. Fine-tuning was performed with Unsloth~\cite{meth_unsloth2023}. BGE-M3 served as the dense retrieval encoder, and BGE-reranker-v2-m3 was used to reorder retrieved passages. These retrieval components were kept fixed when the generators were fine-tuned.

% Figure content: Overview of the four experiments. Show prompting with optional demonstrations; training pairs leading to a fine-tuned generator; a knowledge base feeding the retrieval pipeline and pretrained generator; and the same retrieval pipeline feeding the fine-tuned generator. All four branches should end in an answer to a test question.
\begin{figure}[!t]
\centering
% Replace the box below with \includegraphics[width=\columnwidth]{your_figure_file}.
\fbox{\parbox[c][1.15in][c]{0.94\columnwidth}{\centering Figure placeholder}}
\caption{Experimental configurations for prompting, fine-tuning, RAG, and RAG combined with fine-tuning.}
\label{fig:experiments}
\end{figure}

\subsection{Experiment 1: Prompting}
\label{exp:1}

This experiment evaluated the instruction-tuned models without additional training or external retrieval. We used zero-shot prompting, few-shot prompting, and few-shot chain-of-thought prompting. The prompts followed each model's conversational format and requested answers in Bangla.

\subsubsection{Zero-Shot Prompting}

The model received an instruction and a test question without example answers. The instruction requested a concise response that addressed the information sought in the question, such as a fee, required document, or application procedure. This setting provided a baseline for answering from the model's existing knowledge.

\subsubsection{Few-Shot Prompting}

Few-shot prompting used three example questions and answers~\cite{meth_brown2020fewshot} to establish the expected response format. For each service domain, three pairs were sampled from the training set using random seed 42. After removing duplicate questions and normalized exact matches with test questions, the selected examples were placed before each test question from the corresponding domain and remained fixed throughout the experiment. Qwen, Llama, and Mistral received the same examples in the same order, formatted according to their respective conversational templates. This arrangement provided a common set of demonstrations for comparing the three models.

\subsubsection{Few-Shot Chain-of-Thought Prompting}

The chain-of-thought setting~\cite{meth_wei2022cot} extended the same three-shot procedure by including a brief explanatory rationale before the answer in each demonstration. The selected questions, reference answers, and example order were shared across the three models. Each model was instructed to provide a brief analysis followed by a concise final answer in Bangla. The final response was used for scoring.

Llama and Mistral prompting used greedy decoding, a repetition penalty of 1.05, and a maximum of 500 generated tokens. Qwen also used greedy decoding, beginning with an output budget of 512 tokens. Responses that reached the limit were regenerated with larger budgets, up to 4,096 tokens for zero-shot prompting and 1,024 tokens for the few-shot settings.

\subsection{Experiment 2: Fine-Tuning}
\label{exp:2}

The second experiment adapted each generator to the training QA pairs and evaluated it without retrieved passages. Training used the same 938 examples across the three models. The purpose was to assess the effect of learning the terminology and answer patterns of the selected government services.

% Figure content: Fine-tuning process. Show training QA pairs, conversation formatting, a frozen quantized generator with trainable LoRA adapters, optimization, and evaluation on the test questions.
\begin{figure}[!t]
\centering
% Replace the box below with \includegraphics[width=\columnwidth]{your_figure_file}.
\fbox{\parbox[c][1.15in][c]{0.94\columnwidth}{\centering Figure placeholder}}
\caption{QLoRA fine-tuning of the answer generators.}
\label{fig:qlora_training}
\end{figure}

\subsubsection{Training Data Representation}

Each training example was formatted as a conversation containing a system instruction, a user question, and the reference answer. Additional input accompanied the question where available. The conversation followed the corresponding model's instruction format, and training used an autoregressive language-modeling objective. Sequences were packed with a maximum training length of 1,024 tokens.

\subsubsection{Model Configuration}

We applied quantized low-rank adaptation (QLoRA)~\cite{meth_dettmers2023qlora} to train the generators within the available GPU resources. The backbone weights remained frozen in 4-bit form. Trainable adapters~\cite{meth_hu2022lora} were added to the query, key, value, and output projections of the attention layers, as well as the gate, up, and down projections of the feed-forward layers. The adapter rank and scaling parameter were both set to 16, giving an alpha-to-rank ratio of one. Adapter dropout was zero, and no bias parameters were trained.

\subsubsection{Model Training}

Each model was trained for three epochs using 8-bit AdamW~\cite{meth_loshchilov2019adamw} with a learning rate of 0.0002. A per-device batch size of two and four gradient accumulation steps gave an effective batch size of eight sequences on one GPU. The learning rate followed a cosine schedule, with a warmup ratio of 0.03 and weight decay of 0.01. Unsloth gradient checkpointing was enabled, and the training seed was 3407. Checkpoints were retained after each epoch. Table~\ref{tab:finetuning_settings} summarizes the shared settings.

\begin{table}[!t]
\centering
\caption{Fine-tuning configuration used for the three generators.}
\label{tab:finetuning_settings}
\footnotesize
\setlength{\tabcolsep}{3pt}
\renewcommand{\arraystretch}{1.1}
\begin{tabular}{p{0.43\columnwidth}p{0.49\columnwidth}}
\hline
\textbf{Parameter} & \textbf{Setting} \\
\hline
Adaptation & QLoRA, frozen 4-bit backbone \\
LoRA rank / alpha & 16 / 16 \\
LoRA dropout / bias & 0 / none \\
Epochs & 3 \\
Optimizer & 8-bit AdamW \\
Learning rate & 0.0002 \\
Per-device batch size & 2 \\
Gradient accumulation & 4 steps \\
Effective batch size & 8 sequences \\
Learning-rate schedule & Cosine \\
Warmup ratio & 0.03 \\
Weight decay & 0.01 \\
Maximum sequence length & 1,024 tokens \\
Sequence packing & Enabled \\
Gradient checkpointing & Unsloth \\
Random seed & 3407 \\
\hline
\end{tabular}
\end{table}

The fine-tuned generators then answered the test questions without access to the retrieval index. The output limit was 500 tokens for all three models.

\subsection{Experiment 3: RAG}
\label{exp:3}

The third experiment supplied external evidence to each pretrained generator without domain-specific fine-tuning. The retrieval pipeline combined dense and sparse search, fused their rankings, and reranked the resulting candidates before constructing the generation prompt.

% Figure content: Hybrid retrieval process. Show a question and supplied domain label entering BGE-M3 dense retrieval and BM25 sparse retrieval in parallel. Each contributes up to 100 passages. Show reciprocal rank fusion, the top 25 candidates, BGE reranking, the top five passages, and answer generation.
\begin{figure}[!t]
\centering
% Replace the box below with \includegraphics[width=\columnwidth]{your_figure_file}.
\fbox{\parbox[c][1.15in][c]{0.94\columnwidth}{\centering Figure placeholder}}
\caption{Hybrid retrieval and reranking for government service question answering.}
\label{fig:hybrid_retrieval}
\end{figure}

\subsubsection{Knowledge Base Preparation and Indexing}

Knowledge base records were divided into passages of at most 180 whitespace-delimited words, with an overlap of 30 words between adjacent passages. This produced 1,284 passages across the four domains. The indexed representation incorporated the title, topic, and passage content. Document identifiers, domain labels, and source URLs were retained as metadata.

Dense representations were generated with BGE-M3 (BAAI/bge-m3)~\cite{meth_chen2024bgem3}, using batches of eight passages. Both question and passage embeddings were normalized to unit length. The passage vectors were stored in a FAISS IndexFlatIP index~\cite{meth_johnson2019faiss}, so exact inner-product search produced the same ordering as cosine similarity for these normalized vectors.

A separate sparse index was constructed with BM25Okapi~\cite{meth_robertson2009bm25} using its default parameters. Preprocessing applied Unicode NFKC normalization, converted Bangla digits to Latin digits, and lowercased Latin characters. Tokenization preserved Bangla character sequences, English words, and numerical expressions.

\subsubsection{Hybrid Retrieval and Reranking}

Each question was accompanied by its service-domain label. Dense and sparse retrieval each retained up to 100 passages from that domain. The evaluation therefore assumed that the relevant service category was already known.

The two rankings were combined through reciprocal rank fusion~\cite{meth_cormack2009rrf} with equal weight. For each list containing a passage, its contribution was one divided by the sum of 60 and its rank, with ranks starting at one. Contributions were added across the two lists. The 25 passages with the highest fused scores proceeded to reranking.

BGE-reranker-v2-m3 (BAAI/bge-reranker-v2-m3)~\cite{baai2024reranker} scored each question and candidate passage jointly. Reranking used batches of four and a maximum pair length of 512 tokens. The five highest-scoring passages became the evidence supplied to the generator. Table~\ref{tab:retrieval_settings} lists the retrieval settings.

\begin{table}[!t]
\centering
\caption{Retrieval and reranking configuration.}
\label{tab:retrieval_settings}
\footnotesize
\setlength{\tabcolsep}{3pt}
\renewcommand{\arraystretch}{1.1}
\begin{tabular}{p{0.43\columnwidth}p{0.49\columnwidth}}
\hline
\textbf{Component} & \textbf{Setting} \\
\hline
Passage length / overlap & 180 / 30 words \\
Dense encoder & BAAI/bge-m3 \\
Embedding batch size & 8 \\
Dense index & FAISS IndexFlatIP \\
Sparse retriever & BM25Okapi, default parameters \\
Domain restriction & Supplied question domain \\
Candidates per retriever & Up to 100 \\
Rank fusion & RRF, constant 60, equal weights \\
Candidates for reranking & 25 \\
Reranker & BAAI/bge-reranker-v2-m3 \\
Reranker batch size & 4 \\
Maximum reranker pair length & 512 tokens \\
Passages supplied to generator & 5 \\
\hline
\end{tabular}
\end{table}

\subsubsection{Answer Generation}

The prompt combined the question with the titles and contents of the selected passages. It requested a concise Bangla answer based on that evidence and called attention to fees, processing times, numerical values, and required documents. If the retrieved material did not contain the requested information, the model was instructed to acknowledge this.

The generators received the same selected passages for each question. RAG inference used greedy decoding, a repetition penalty of 1.05, and an input limit of 7,000 tokens. The output limit 500 tokens for three models.

\subsection{Experiment 4: Combining Fine-Tuning and RAG}
\label{exp:4}

The final experiment combined the generators adapted in Experiment~\ref{exp:2} with the retrieval pipeline from Experiment~\ref{exp:3}. Each fine-tuned model received a test question and the same five passages supplied to its pretrained counterpart. The knowledge base, dense encoder, BM25 index, fusion procedure, and reranker were unchanged. No additional training of the retrieval components was performed.

This comparison examined whether adaptation on the government service QA pairs helped a generator answer when external evidence was available. Comparing the combined setting with standalone fine-tuning assessed the contribution of retrieval. Comparing it with RAG assessed the effect of generator adaptation under the same evidence. The combined setting used greedy decoding, a repetition penalty of 1.05, a maximum input length of 7,000 tokens, and an output limit of 500 tokens for all three generators.

\section{Results and Discussion}

This section presents the experimental results of the evaluated large
language models for Bangla government-service question answering. The
analysis covers prompting-based baselines, supervised fine-tuning,
Retrieval-Augmented Generation (RAG), and the hybrid RAG + Fine-Tuning
configuration. Performance is examined using exact, lexical, and semantic
evaluation metrics, followed by analysis of adaptation strategies and
cross-model behavior.

\subsection{Evaluation Metrics}

The evaluated models were assessed using multiple automatic metrics to
capture exact agreement, lexical overlap, and semantic similarity between
generated responses and reference answers. Since no single metric fully
represents generative answer quality, the metrics were considered jointly.

\textbf{Exact Match (EM)} measures the proportion of generated responses
that exactly match the reference answers after normalization.

\textbf{Fuzzy Match} measures approximate string similarity and therefore
allows partial agreement when the generated and reference responses differ
slightly in wording.

\textbf{BLEU} evaluates n-gram precision between generated and reference
responses, while \textbf{ROUGE} measures recall-oriented lexical overlap.
ROUGE-1, ROUGE-2, and ROUGE-L were used to evaluate unigram, bigram, and
longest-common-subsequence similarity, respectively.

\textbf{Token F1} computes the harmonic mean of token-level precision and
recall and therefore provides partial credit when only part of the reference
answer is reproduced.

\textbf{BERTScore} measures semantic similarity using contextualized token
representations. It is useful when generated and reference answers convey
similar meaning using different lexical forms.

Together, these metrics provide complementary views of performance. EM
reflects exact correspondence, BLEU and ROUGE capture lexical overlap,
Token F1 measures partial agreement, and BERTScore captures semantic
similarity. Consequently, the reported results are interpreted collectively
rather than relying on any single metric.

\subsection{Overall Performance Comparison}

Table~\ref{tab:overall_results} summarizes the performance of all evaluated
configurations. A clear progression is observed from prompting-based
baselines toward retrieval-based and hybrid approaches. Zero-shot and
few-shot settings generally produced the lowest scores, whereas RAG
substantially improved performance across all three model families. The
strongest overall results were obtained when retrieval was combined with
fine-tuning.

\begin{table*}[t]
\centering
\caption{Performance Comparison of Different LLM Configurations for Bangla
Government-Service Question Answering}
\label{tab:overall_results}
\resizebox{\textwidth}{!}{
\begin{tabular}{lcccccc}
\hline
\textbf{Model / Configuration} &
\textbf{EM} &
\textbf{Fuzzy} &
\textbf{BLEU} &
\textbf{ROUGE-L} &
\textbf{Token F1} &
\textbf{BERT F1} \\
\hline
Llama Zero-Shot           & 0.000 & 0.513 & 0.018 & 0.116 & 0.137 & 0.658 \\
Llama Few-Shot            & 0.000 & 0.576 & 0.050 & 0.238 & 0.262 & 0.726 \\
Llama Fine-Tuned          & 0.004 & 0.602 & 0.134 & 0.286 & 0.320 & 0.779 \\
Llama RAG                 & 0.060 & 0.730 & 0.338 & 0.433 & 0.449 & 0.797 \\
Llama RAG + Fine-Tuning   & 0.355 & 0.908 & 0.478 & 0.673 & 0.687 & 0.888 \\
\hline
Qwen Zero-Shot            & 0.000 & 0.497 & 0.016 & 0.139 & 0.168 & 0.694 \\
Qwen Few-Shot             & 0.000 & 0.511 & 0.027 & 0.177 & 0.208 & 0.715 \\
Qwen Few-Shot + CoT       & 0.000 & 0.527 & 0.042 & 0.205 & 0.233 & 0.727 \\
Qwen RAG                  & 0.169 & 0.881 & 0.412 & 0.622 & 0.654 & 0.856 \\
Qwen RAG + Fine-Tuning    & \textbf{0.391} & \textbf{0.924} &
\textbf{0.636} & \textbf{0.771} & \textbf{0.788} &
\textbf{0.923} \\
\hline
Mistral Zero-Shot         & 0.000 & 0.491 & 0.013 & 0.106 & 0.127 & 0.662 \\
Mistral Few-Shot          & 0.000 & 0.482 & 0.013 & 0.122 & 0.146 & 0.670 \\
Mistral Few-Shot + CoT    & 0.000 & 0.494 & 0.022 & 0.143 & 0.167 & 0.694 \\
Mistral RAG               & 0.008 & 0.762 & 0.221 & 0.422 & 0.454 & 0.873 \\
Mistral RAG + Fine-Tuning & 0.169 & 0.868 & 0.506 & 0.630 & 0.652 & 0.888 \\
\hline
\end{tabular}}
\end{table*}

At the baseline level, all three zero-shot models obtained an EM score of
0.000, indicating limited ability to reproduce exact reference answers using
pretrained knowledge alone. Few-shot and CoT prompting produced only modest
improvements. In contrast, RAG resulted in substantial gains, and the hybrid
configuration further improved performance.

Among all evaluated configurations, Qwen RAG + Fine-Tuning achieved the
highest reported scores, including 0.391 EM, 0.924 Fuzzy Match, 0.636 BLEU,
0.771 ROUGE-L, 0.788 Token F1, and 0.923 BERT F1. Llama and Mistral also
benefited considerably from the hybrid configuration, indicating that the
combination of retrieval and domain adaptation consistently improves answer
quality.

To complement the numerical results in Table~\ref{tab:overall_results},
Fig.~\ref{fig:multimetric_comparison} visualizes four representative
metrics---EM, ROUGE-L, Token F1, and BERT F1---for the Zero-shot, RAG,
and RAG + Fine-Tuning configurations. The figure highlights the progression
from prompting-only generation to retrieval-grounded and hybrid approaches
for all three model families.

\begin{figure*}[t]
    \centering
    \includegraphics[width=\textwidth]
    {comp.png}
    \caption{Multi-metric comparison of representative Llama, Qwen, and
    Mistral configurations using EM, ROUGE-L, Token F1, and BERT F1.
    The figure illustrates the performance progression from Zero-shot to
    RAG and RAG + Fine-Tuning.}
    \label{fig:multimetric_comparison}
\end{figure*}

The figure also illustrates that the different evaluation metrics do not
always move in parallel. For example, some configurations achieve relatively
high semantic similarity despite low Exact Match performance. This behavior
is particularly visible for Mistral RAG, which obtains a BERT F1 score of
0.873 but an EM score of only 0.008. Such differences motivate the use of
multiple complementary metrics instead of relying on a single performance
measure.

\subsection{Effect of Prompting, Fine-Tuning, and RAG}

The evaluated adaptation strategies affected model performance in different
ways. Prompting-based methods provide task demonstrations or reasoning
instructions without modifying the model parameters or introducing new
external knowledge. Consequently, zero-shot, few-shot, and CoT configurations
showed comparatively limited gains, particularly for questions requiring
specific service-related facts.

Fine-tuning improved domain adaptation by exposing the model to task-specific
question--answer pairs. For Llama, Token F1 increased from 0.262 under
few-shot prompting to 0.320 after fine-tuning, while BERT F1 increased from
0.726 to 0.779. BLEU also increased from 0.050 to 0.134. These gains suggest
that fine-tuning helps the model learn domain terminology, expected response
patterns, and task-specific linguistic characteristics.

However, the fine-tuned Llama obtained an EM score of only 0.004. This
indicates that learning domain patterns through parameter updates alone does
not guarantee precise reproduction of service-specific facts. Government
service questions often require exact information related to fees, required
documents, eligibility conditions, processing details, and procedural steps.

RAG produced a more pronounced improvement because it supplied relevant
government-service information during inference. Llama RAG increased Token
F1 to 0.449 and ROUGE-L to 0.433. Qwen RAG achieved 0.169 EM, 0.881
Fuzzy Match, 0.622 ROUGE-L, 0.654 Token F1, and 0.856 BERT F1. Mistral
RAG similarly improved to 0.454 Token F1 and 0.873 BERT F1. These results
show that access to external domain knowledge provides a clear advantage
over prompting-only configurations.

The hybrid RAG + Fine-Tuning configuration combines the advantages of both
adaptation mechanisms. Fine-tuning helps the model learn domain-specific
language and response structure, whereas retrieval supplies relevant factual
context during generation. This complementary effect is reflected in the
performance improvements observed across all three model families.

For Llama, the hybrid configuration reached 0.355 EM, 0.673 ROUGE-L,
0.687 Token F1, and 0.888 BERT F1. Mistral improved to 0.169 EM,
0.630 ROUGE-L, 0.652 Token F1, and 0.888 BERT F1. Qwen produced the
strongest hybrid result, with substantial gains over its RAG-only
configuration. These results indicate that retrieval and fine-tuning address
different aspects of the task and are more effective when applied together.

\subsection{Cross-Model Comparison}

The three model families also responded differently to the evaluated
adaptation strategies. Under prompting-only settings, their performance
remained relatively close, and no large separation was observed among the
models. More noticeable differences emerged after retrieval was introduced.

Qwen benefited strongly from retrieval and achieved the highest scores across
the common metrics in the hybrid configuration. Llama also exhibited
substantial improvements after fine-tuning and retrieval, particularly in EM,
ROUGE-L, Token F1, and BERT F1. Mistral showed strong semantic similarity in
retrieval-based settings, although its exact-match performance remained
lower.

The contrast between Mistral RAG's BERT F1 of 0.873 and EM of 0.008 is a
clear example of how model evaluation depends on the metric being considered.
A response may be semantically similar to a reference answer while differing
in exact wording, numerical values, or specific factual details. Conversely,
a strong lexical match does not necessarily guarantee that all facts in the
response are correct.

These observations are particularly important in government-service QA,
where seemingly small factual differences may change the meaning of an
answer. Incorrect fees, eligibility conditions, required documents, deadlines,
or procedural steps can make an otherwise fluent response unreliable for
practical use.

\subsection{Implications of Evaluation Metrics and Source Verification}

The experimental results expose an important limitation of conventional
automatic evaluation metrics. EM, BLEU, ROUGE, Token F1, Fuzzy Match, and
BERTScore capture different aspects of similarity between generated and
reference answers, but none of these metrics directly guarantees factual
correctness.

Exact Match is highly restrictive and can assign a score of zero even when
a generated response conveys relevant information using different wording.
In contrast, semantic metrics such as BERTScore may assign a relatively high
similarity score even when a response contains an incorrect number, fee,
document requirement, or procedural condition. Lexical overlap metrics
occupy an intermediate position but also remain dependent on the wording of
the reference answer.

For government-service applications, these limitations are particularly
important because response reliability depends not only on linguistic
similarity but also on factual consistency with authoritative information.
Government policies, fees, procedural requirements, and service instructions
may also change over time, making temporal freshness an additional concern.

Therefore, the findings support evaluating generated responses using multiple
automatic metrics while complementing them with source-based verification.
Fact-sensitive responses should be checked against reliable and up-to-date
government sources rather than being considered correct solely because they
achieve high lexical or semantic similarity scores. This source-grounded
perspective is especially relevant for RAG-based government-service systems,
where retrieved evidence can provide a direct basis for validating generated
answers.

Overall, the experimental findings show that retrieval-based grounding is
critical for improving Bangla government-service question answering, while
fine-tuning provides additional benefit by adapting models to the target
domain. Their combination delivers the strongest overall performance across
the evaluated models. At the same time, the observed differences among
automatic metrics demonstrate that quantitative similarity scores alone are
insufficient for establishing answer reliability, reinforcing the importance
of verification against authoritative and current government information.

\section{Limitations, Reliability, and Deployment Considerations}
\label{sec:limitations}

Evaluation was confined to four Bangladesh government-service domains: passport and e-passport, National Identification (NID), birth and death registration, and Taxpayer Identification Number (TIN). The dataset contains prepared question--answer pairs derived from the constructed knowledge base rather than naturally collected citizen conversations. Its held-out questions therefore assess answering over the collected information. They may not capture the range of ambiguous or incomplete requests, spelling variation, conversational context, and unexpected phrasing encountered in practical use.

Paraphrased questions were generated from the same service information. Questions associated with related knowledge-base records may consequently remain semantically similar across the training and test sets, even when exact duplicates are excluded. This possibility does not establish that data leakage occurred, but it limits what can be inferred about generalization to unfamiliar intents. A stricter evaluation could group examples by the underlying record, information item, or intent before partitioning the dataset.

Retrieval used the service-domain label supplied with each test question. This restriction kept unrelated service categories out of candidate selection and provided a controlled comparison. A deployed system would ordinarily need to identify the domain from the question or search across the complete knowledge base. Errors at this stage could affect the evidence available to the generator.

Retrieval ranking was evaluated using Hit@1, Hit@3, Hit@5, MRR@5, and nDCG@5. The relevance judgments were derived from reference answers rather than a fully manually annotated passage dataset. These measures provide a useful indication of ranking quality under the adopted relevance definition. Manually annotated passage relevance would provide stronger evidence in future work.

Government information can change after collection, including fees, eligibility rules, required documents, procedures, processing times, and online links. RAG permits updates to the external knowledge base without retraining the generator, but does not guarantee factual correctness. Reliability still depends on current source information, retrieval of the relevant evidence, and generation that neither omits nor misinterprets that evidence or introduces unsupported details. Source URLs and passage metadata support traceability, provided that the collection and its index are maintained.

Exact Match, Fuzzy Match, BLEU, ROUGE, Token F1, and BERTScore measure exact, lexical, or semantic agreement with reference answers. A strong score can coexist with an incorrect fee, document requirement, eligibility condition, processing time, or procedural detail. These measures therefore cannot establish factual reliability on their own.

The approaches also impose different operational requirements. Prompting needs neither generator adaptation nor an external retrieval index. Fine-tuning requires adapting the generator, although QLoRA reduces the memory needed for this process. RAG adds an indexed knowledge base, retrieval, rank fusion, reranking, and evidence construction during inference. Combining RAG with fine-tuning retains both adaptation and retrieval infrastructure, increasing operational complexity relative to prompting or standalone fine-tuning. Citizen-facing use would require regular knowledge-base updates, automatic domain identification or unrestricted retrieval, accessible source information, and a response policy for questions lacking sufficient evidence. Fact-sensitive details would need verification before generated answers were relied upon.

The experiments establish performance within the evaluated Bangla government-service setting. Broader deployment requires assessment on more diverse natural queries, continuously updated authoritative information, and stronger factual validation.

\section{Conclusion and Future Work}
\label{sec:conclusion}
This study investigated Bangla government-service question answering
across three open-weight language model families---Qwen, Llama, and
Mistral---using prompting, supervised fine-tuning, Retrieval-Augmented
Generation (RAG), and RAG combined with fine-tuning. A domain-specific
dataset containing 1,186 question--answer pairs was constructed from
official government information covering birth and death registration,
passport and e-passport services, National Identification (NID), and
Taxpayer Identification Number (TIN) services. A corresponding knowledge
base was prepared to support retrieval-grounded generation.

The experimental results show a clear advantage for retrieval-based
approaches over prompting-only configurations. Zero-shot and few-shot
methods produced comparatively low performance, while RAG substantially
improved lexical and semantic agreement with the reference answers across
all three model families. Fine-tuning improved domain-specific response
patterns, but its effect was more limited when external factual information
was unavailable. The strongest overall performance was obtained by
combining retrieval with fine-tuning, indicating that the two approaches
address complementary aspects of the task: fine-tuning adapts the
generator to the target domain, whereas retrieval supplies relevant
government-service information at inference time.

Among the evaluated configurations, Qwen with RAG and fine-tuning achieved
the best overall result, reaching an Exact Match of 0.391, Fuzzy Match of
0.924, BLEU of 0.636, ROUGE-L of 0.771, Token F1 of 0.788, and BERT F1 of
0.923. Llama and Mistral also showed substantial improvements under the
hybrid configuration, supporting the general benefit of combining
retrieval grounding with domain adaptation.

At the same time, the results demonstrate that automatic similarity
metrics alone are insufficient to establish the factual reliability of a
government-service answer. A response may obtain a high semantic or
lexical similarity score while still containing an incorrect fee,
eligibility condition, document requirement, or procedural detail.
Therefore, practical government-service systems should evaluate generated
answers together with the retrieved evidence and verify fact-sensitive
information against authoritative and current government sources.

The present study is limited to four government-service domains and a
dataset generated from the prepared government knowledge base rather than
naturally occurring citizen interactions. In addition, the retrieval
experiments assume that the service domain of each question is already
known. Future work should therefore evaluate the system on real citizen
queries, extend the knowledge base to additional public services, and
introduce automatic service-domain identification. Further work should
also examine retrieval accuracy directly, incorporate source-based factual
verification, and study mechanisms for detecting unsupported or outdated
answers. Human evaluation involving Bangla-speaking users or domain
experts would provide an additional measure of factual correctness,
clarity, and practical usefulness beyond automatic text-similarity
metrics.

Overall, the findings indicate that retrieval-grounded generation combined
with domain-specific fine-tuning is the most effective strategy among the
evaluated approaches for Bangla government-service question answering.
The results also emphasize that strong generation performance should be
accompanied by reliable retrieval, up-to-date authoritative information,
and explicit factual verification before such systems are used in
citizen-facing public-service applications.


% Intentionally left blank.

\bibliographystyle{IEEEtran}
\bibliography{reference}

\end{document}
